\documentclass{article}

% -----------------------------
% Packages
% -----------------------------
\usepackage{fancyhdr}
\usepackage{extramarks}
\usepackage{amsmath}
\usepackage{amsthm}
\usepackage{amsfonts}
\usepackage{tikz}
\usepackage[plain]{algorithm}
\usepackage{algpseudocode}
\usepackage{amssymb}
\usepackage{enumitem}

\usetikzlibrary{automata,positioning}

% -----------------------------
% Basic Document Settings
% -----------------------------
\topmargin=-0.45in
\evensidemargin=0in
\oddsidemargin=0in
\textwidth=6.5in
\textheight=9.0in
\headsep=0.25in

\linespread{1.1}

\pagestyle{fancy}
\lhead{\hmwkAuthorName}
\chead{\hmwkClass: \hmwkTitle}
\rhead{}
\lfoot{\lastxmark}
\cfoot{\thepage}

\renewcommand\headrulewidth{0.4pt}
\renewcommand\footrulewidth{0.4pt}

\setlength\parindent{0pt}

% -----------------------------
% Problem Section Helpers
% -----------------------------
\newcommand{\enterProblemHeader}[1]{
  \nobreak\extramarks{}{Problem #1 continued on next page\ldots}\nobreak{}
  \nobreak\extramarks{Problem #1 (continued)}{Problem #1 continued on next page\ldots}\nobreak{}
}

\newcommand{\exitProblemHeader}[1]{
  \nobreak\extramarks{Problem #1 (continued)}{Problem #1 continued on next page\ldots}\nobreak{}
  \nobreak\extramarks{Problem #1}{}\nobreak{}
}

\setcounter{secnumdepth}{0}
\newcounter{partCounter}
\newcounter{homeworkProblemCounter}
\setcounter{homeworkProblemCounter}{1}
\nobreak\extramarks{Problem \arabic{homeworkProblemCounter}}{}\nobreak{}

% -----------------------------
% Theorem env (optional)
% -----------------------------
\newtheorem*{theorem}{Theorem}

% -----------------------------
% Homework Problem Environment
%   Usage:
%   \begin{homeworkProblem}            % auto-number
%     % Your answer here
%   \end{homeworkProblem}
%
%   \begin{homeworkProblem}[1.2.3]     % label problem "1.2.3" (e.g. textbook numbering)
%     % Your answer here
%   \end{homeworkProblem}
% -----------------------------
\newenvironment{homeworkProblem}[1][]{
  \def\hwProblemLabel{#1}
  \ifx\hwProblemLabel\empty
    \edef\hwProblemLabel{\arabic{homeworkProblemCounter}}
    \stepcounter{homeworkProblemCounter}
  \fi
  \section{Problem \hwProblemLabel}
  \setcounter{partCounter}{1}
  \enterProblemHeader{\hwProblemLabel}
}{
  \exitProblemHeader{\hwProblemLabel}
}

% -----------------------------
% Homework Details (EDIT THESE)
% -----------------------------
\newcommand{\hmwkTitle}{Homework \#7}
\newcommand{\hmwkClass}{Math 3320}
\newcommand{\hmwkAuthorName}{\textbf{Joseph Quinn}}

% -----------------------------
% Title
% -----------------------------
\title{
  \vspace{2in}
  \textmd{\textbf{\hmwkClass:\ \hmwkTitle}}\\
  \vspace{3in}
}
\author{\hmwkAuthorName}
\date{}

% -----------------------------
% Convenience Macros (optional)
% -----------------------------
\renewcommand{\part}[1]{\textbf{\large Part \Alph{partCounter}}\stepcounter{partCounter}\\}

% Algorithms
\newcommand{\alg}[1]{\textsc{\bfseries \footnotesize #1}}

% Calculus
\newcommand{\deriv}[1]{\frac{\mathrm{d}}{\mathrm{d}x} (#1)}
\newcommand{\pderiv}[2]{\frac{\partial}{\partial #1} (#2)}
\newcommand{\dx}{\mathrm{d}x}

% Statistics
\newcommand{\E}{\mathrm{E}}
\newcommand{\Var}{\mathrm{Var}}
\newcommand{\Cov}{\mathrm{Cov}}
\newcommand{\Bias}{\mathrm{Bias}}

% Proof step (for aligned derivations)
\newcommand{\step}[2]{& #1 & & \text{#2} \\}

% Solution header (optional)
\newcommand{\solution}{\textbf{\large Solution}}

% -----------------------------
% Document
% -----------------------------
\begin{document}

\begin{homeworkProblem}[A14]

	Given the LP
	\[
		\begin{aligned}
			\max        & \quad z = -6x_1 - 25x_2 - 14x_3 \\
			\text{s.t.} & -x_1 - 2x_2 - 4x_3  \leq -46    \\
			            & 4x_1 - 5x_2 - 6x_3  \leq -54    \\
			            & -x_1 - 4x_2 - 2x_3  \leq -24    \\
			            & x_1, x_2, x_3       \geq 0
		\end{aligned}
	\]

	\begin{enumerate}
		\item[(I2)] Adding slack variables $x_4, x_5, x_6$ gives the initial dictionary,
		      \[
			      \begin{aligned}
				      x_4 & = -46 + x_1 + 2x_2 + 4x_3  \\
				      x_5 & = -54 - 4x_1 + 5x_2 + 6x_3 \\
				      x_6 & = -24 + x_1 + 4x_2 + 2x_3  \\[4pt]
				      z   & = -6x_1 - 25x_2 - 14x_3
			      \end{aligned}
		      \]
		\item[(I3)] Every coefficient in $z$ is already $\le 0$, so this dictionary is dual feasible. But $x_4,x_5,x_6$ are all negative, making it primal infeasible, so we can apply the dual simplex method.
	\end{enumerate}

	\begin{enumerate}
		\item[(D1)] The constants are $-46$, $-54$, and $-24$; by the dual largest coefficient rule we take the one of largest magnitude, so $x_5=-54$ is our leaving variable.
		\item[(D2)] Looking at $x_5$'s row, only $x_2$ and $x_3$ are positive meaning the dual ratio test gives,
		      \[
			      \frac{|{-25}|}{5} = 5 \text{ for } x_2, \qquad \frac{|{-14}|}{6} = \frac{7}{3} \text{ for } x_3,
		      \]
		      so $x_3$ has the smaller ratio and becomes our entering variable.
		\item[(D3)] After performing the pivot we get the dictionary,
		      \[
			      \begin{aligned}
				      x_3 & = 9 + \tfrac{2}{3}x_1 - \tfrac{5}{6}x_2 + \tfrac{1}{6}x_5      \\
				      x_4 & = -10 + \tfrac{11}{3}x_1 - \tfrac{4}{3}x_2 + \tfrac{2}{3}x_5   \\
				      x_6 & = -6 + \tfrac{7}{3}x_1 + \tfrac{7}{3}x_2 + \tfrac{1}{3}x_5     \\[4pt]
				      z   & = -126 - \tfrac{46}{3}x_1 - \tfrac{40}{3}x_2 - \tfrac{7}{3}x_5
			      \end{aligned}
		      \]
	\end{enumerate}

	\begin{enumerate}
		\item[(D1)] From the rule $x_4$ has the larger magnitude, so it is our leaving variable.
		\item[(D2)] Looking at $x_4$'s row, only $x_1$ and $x_5$ are positive meaning the dual ratio test gives,
		      \[
			      \frac{|{-46/3}|}{11/3} = \frac{46}{11} \text{ for } x_1, \qquad \frac{|{-7/3}|}{2/3} = \frac{7}{2} \text{ for } x_5,
		      \]
		      so $x_5$ has the smaller ratio and becomes our entering variable.
		\item[(D3)] After performing the pivot we get the dictionary,
		      \[
			      \begin{aligned}
				      x_3 & = \tfrac{23}{2} - \tfrac14 x_1 - \tfrac12 x_2 + \tfrac14 x_4 \\
				      x_5 & = 15 - \tfrac{11}{2}x_1 + 2x_2 + \tfrac32 x_4                \\
				      x_6 & = -1 + \tfrac12 x_1 + 3x_2 + \tfrac12 x_4                    \\[4pt]
				      z   & = -161 - \tfrac52 x_1 - 18x_2 - \tfrac72 x_4
			      \end{aligned}
		      \]
	\end{enumerate}

	\begin{enumerate}
		\item[(D1)] Since $x_6=-1$ is the only negative constant, it is our leaving variable.
		\item[(D2)] Looking at $x_6$'s row, $x_1$, $x_2$, and $x_4$ all have positive coefficients, so they are all eligible. The dual ratio test gives,
		      \[
			      \frac{|{-5/2}|}{1/2} = 5 \text{ for } x_1, \qquad \frac{|{-18}|}{3} = 6 \text{ for } x_2, \qquad \frac{|{-7/2}|}{1/2} = 7 \text{ for } x_4,
		      \]
		      so $x_1$ has the smallest ratio and becomes our entering variable.
		\item[(D3)] After performing the pivot we get the dictionary,
		      \[
			      \begin{aligned}
				      x_1 & = 2 - 6x_2 - x_4 + 2x_6                  \\
				      x_3 & = 11 + x_2 + \tfrac12 x_4 - \tfrac12 x_6 \\
				      x_5 & = 4 + 35x_2 + 7x_4 - 11x_6               \\[4pt]
				      z   & = -166 - 3x_2 - x_4 - 5x_6
			      \end{aligned}
		      \]

		      All basic variable constants are now nonnegative and every coefficient in the $z$ row is still negative, so this dictionary is both primal and dual feasible, meaning it is optimal. The optimal solution is $(x_1=2,x_2=0,x_3=11)$, giving $z=-166$.
	\end{enumerate}

\end{homeworkProblem}

\begin{homeworkProblem}[A15]
	Given the LP
	\[
		\begin{aligned}
			\max \quad z  & = 70x_1 + 250x_2 \\[4pt]
			\text{s.t.} \quad
			2x_1 + 1.5x_2 & \leq 120         \\
			6x_1 + 20x_2  & \leq 600         \\
			3x_1 + 1.6x_2 & \leq 100         \\
			x_1, x_2      & \geq 0
		\end{aligned}
	\]

	With the optimal dictionary
	\[
		\begin{aligned}
			x_3 & = 75 - 1.55x_1 + 0.075x_4 \\
			x_2 & = 30 - 0.3x_1 - 0.05x_4   \\
			x_5 & = 52 - 2.52x_1 + 0.08x_4  \\[4pt]
			z   & = 7500 - 5x_1 - 12.5x_4
		\end{aligned}
	\]

	\begin{enumerate}[label=(\alph*)]
		\item To see how much profit the filing cabinets would have to improve before we include them in our optimal solution, we need to look at its
		      coefficient in the $z$ equation.

		      We know it's a nonbasic variable and thus has a value of 0 and is not produced at all in our optimal solution, the coefficient on $x_1$ tells
		      us how much it would need to increase before it would be included as a basic variable in the optimal solution. In this case
		      it would have to increase by 5 or more.

		      We see this by updating the original $z$ equation to the form of $z=(70+t)x_1+250x_2$

		      This causes our optimal dictionary to have its $z$ equation as $7500 +(t-5)x_1-12.5x_4$. And our goal is to make $x_1$ no longer a nonbasic variable
		      through a pivot meaning $t-5 \ge 0$ or $t \ge 5$

		\item To find out the amount the steel supply could drop before the optimal solution changes we need to look at the $x_5$ basic variable, if we look we can see for the optimal solution
		      $x_5=52$. Since the original constraint constant is 100, in our optimal solution our slack variable is 52 meaning we could reduce the supply down until 48 without affecting the optimal solution.
		\item The current price of wood is \$8.40, if we cannot buy any more wood at this price but another supplier is willing to sell us more we need to determine the marginal value of the wood resource to
		      determine how much of a premium we are willing to pay for the wood.

		      We want to look at the nonbasic slack variable $x_4$ which corresponds to the marginal value of wood. To demonstrate this lets set the value of $x_4=-1$ meaning we use an additional unit of wood,
		      this changes our $z$ value to become, $z=7500-12.5(-1) = 7512.5$ This is a 12.5 unit profit increase over the original optimal solution meaning for each additional unit of wood we increase our profit
		      by 12.5 profit units. So we should be willing to pay up to a 12.5 profit unit margin for this resource.
	\end{enumerate}
\end{homeworkProblem}
\end{document}


